%% BPMN — Modelo de proceso de negocio AS-IS / TO-BE (plantilla)
%% Compilar: tectonic -X compile bpmn-proceso.tex
\documentclass[11pt,a4paper]{article}
\usepackage{../../docs/latex/legasoft}
\usepackage{../../docs/latex/legasoft-diagramas}

\lgtitulo{BPMN: modelo de proceso de negocio}
\lgtituloCorto{BPMN — Proceso}
\lgcodigo{LGS-MOD-BPMN-1}
\lgversion{0.1}
\lgestado{Plantilla}
\lgfecha{\campo{fecha}}
\lgautor{\lgDirector\ (\lgDirectorRol)}

\begin{document}

\begin{center}
{\sffamily\bfseries\LARGE\color{lgPrimario} BPMN: modelo de proceso de negocio\par}
\vspace{2pt}
{\sffamily\normalsize\color{lgGris} \campo{nombre del proceso} · Legasoft\par}
\end{center}
\vspace{6pt}
{\color{lgBorde}\hrule height 0.8pt}
\vspace{10pt}

\begin{porcompletar}[Cómo usar esta plantilla]
El diagnóstico de procesos es el punto de partida de toda intervención de
Legasoft. Se modela primero el proceso actual (\textbf{AS-IS}) para hacer
visibles las tareas manuales y las ineficiencias, y después el proceso propuesto
(\textbf{TO-BE}) que incorpora la digitalización. Duplique la sección del
diagrama para representar cada uno.
\end{porcompletar}

% =============================================================================
\section*{Proceso actual (AS-IS)}

\begin{figure}[H]
\centering
\ajustaancho{%
\begin{tikzpicture}[bpmn, x=1cm, y=1cm]

  % --- Carril 1: solicitante -----------------------------------------------
  \node[bpmninicio] (ini) at (0,0) {};
  \node[bpmnetq, below=1pt of ini, fill=none] {\campo{solicitud recibida}};

  \node[bpmntarea, right=1.1cm of ini] (t1) {\campo{Registrar solicitud (manual)}};
  \node[bpmntarea, right=1.1cm of t1]  (t2) {\campo{Verificar datos}};
  \node[bpmngateway, right=1.1cm of t2] (g1) {};
  \node[bpmnetq, above=2pt of g1, fill=none] {\campo{¿completa?}};

  \node[bpmntarea, right=1.4cm of g1] (t3) {\campo{Aprobar y archivar}};
  \node[bpmnfin, right=1.1cm of t3] (fin) {};
  \node[bpmnetq, below=1pt of fin, fill=none] {\campo{proceso cerrado}};

  \node[bpmntarea, below=1.3cm of g1] (t4) {\campo{Solicitar corrección}};

  % --- Flujos ---------------------------------------------------------------
  \draw[bpmnflujo] (ini) -- (t1);
  \draw[bpmnflujo] (t1) -- (t2);
  \draw[bpmnflujo] (t2) -- (g1);
  \draw[bpmnflujo] (g1) -- (t3) node[bpmnetq, pos=0.45] {sí};
  \draw[bpmnflujo] (t3) -- (fin);
  \draw[bpmnflujo] (g1) -- (t4) node[bpmnetq, pos=0.5] {no};
  \draw[bpmnflujo] (t4.west) -- ++(-1.0,0) |- (t2.south);

\end{tikzpicture}}
\caption{Proceso actual. Las tareas marcadas como manuales son las candidatas a
digitalización.}
\end{figure}

% =============================================================================
\Needspace*{20\baselineskip}
\section*{Hallazgos del diagnóstico}

\begin{center}
\begin{tabularx}{\linewidth}{@{}C{1.7cm} Y L{3.4cm}@{}}
\toprule
\sffamily\bfseries\color{lgPrimario}ID &
\sffamily\bfseries\color{lgPrimario}Hallazgo &
\sffamily\bfseries\color{lgPrimario}Consecuencia observada \\
\midrule
H-01 & \campo{tarea manual o repetitiva identificada} & \campo{tiempo, error o retrabajo} \\
H-02 & \campo{registro disperso o duplicado}          & \campo{información no confiable} \\
H-03 & \campo{falta de control o trazabilidad}        & \campo{imposibilidad de auditar} \\
\bottomrule
\end{tabularx}
\captionof{table}{Hallazgos que justifican la propuesta de mejora.}
\end{center}

% =============================================================================
\Needspace*{20\baselineskip}
\section*{Proceso propuesto (TO-BE)}

\begin{figure}[H]
\centering
\ajustaancho{%
\begin{tikzpicture}[bpmn, x=1cm, y=1cm]

  \node[bpmninicio] (ini2) at (0,0) {};
  \node[bpmnetq, below=1pt of ini2, fill=none] {\campo{solicitud en línea}};

  \node[bpmntarea, right=1.1cm of ini2] (u1) {\campo{Registrar en el sistema}};
  \node[bpmntarea, right=1.1cm of u1]  (u2) {\campo{Validación automática}};
  \node[bpmngateway, right=1.1cm of u2] (g2) {};
  \node[bpmnetq, above=2pt of g2, fill=none] {\campo{¿válida?}};

  \node[bpmntarea, right=1.4cm of g2] (u3) {\campo{Aprobar con registro de auditoría}};
  \node[bpmnfin, right=1.1cm of u3] (fin2) {};

  \node[bpmntarea, below=1.3cm of g2] (u4) {\campo{Notificar al solicitante}};

  \draw[bpmnflujo] (ini2) -- (u1);
  \draw[bpmnflujo] (u1) -- (u2);
  \draw[bpmnflujo] (u2) -- (g2);
  \draw[bpmnflujo] (g2) -- (u3) node[bpmnetq, pos=0.45] {sí};
  \draw[bpmnflujo] (u3) -- (fin2);
  \draw[bpmnflujo] (g2) -- (u4) node[bpmnetq, pos=0.5] {no};
  \draw[bpmnmensaje] (u4.west) -- ++(-1.0,0) |- (u1.south);

\end{tikzpicture}}
\caption{Proceso propuesto con validación automática, notificación al solicitante
y registro de auditoría.}
\end{figure}

% =============================================================================
\Needspace*{20\baselineskip}
\section*{Notación empleada}

\begin{center}
\begin{tabularx}{\linewidth}{@{}L{4.2cm} Y@{}}
\toprule
\sffamily\bfseries\color{lgPrimario}Elemento &
\sffamily\bfseries\color{lgPrimario}Significado \\
\midrule
Círculo de borde fino     & Evento de inicio: qué dispara el proceso. \\
Rectángulo redondeado     & Tarea o actividad realizada por un responsable. \\
Rombo                     & Compuerta de decisión; las salidas se rotulan con la condición. \\
Círculo de borde gruesa   & Evento de fin: estado en que termina el proceso. \\
Flecha continua           & Flujo de secuencia dentro del proceso. \\
Flecha discontinua        & Flujo de mensaje entre participantes distintos. \\
\bottomrule
\end{tabularx}
\captionof{table}{Elementos BPMN utilizados en los modelos de Legasoft.}
\end{center}

\begin{nota}[Alcance de la propuesta]
La comparación AS-IS / TO-BE sustenta la propuesta de mejora ante el cliente. La
digitalización debe ser proporcional al problema: si una tarea manual no genera
error ni retrabajo medible, automatizarla no aporta un beneficio verificable.
\end{nota}

\end{document}
